\section*{Literaturverzeichnis}
\addcontentsline{toc}{section}{Literaturverzeichnis}

\begin{thebibliography}{99}

\bibitem{aburbeian2024}
Aburbeian, A. M., \& Fernández-Veiga, M. (2024). Secure Internet Financial Transactions: A Framework Integrating Multi-Factor Authentication and Machine Learning. \textit{AI}, 5(1), 177--194.

\bibitem{acemyan2018}
Acemyan, C. Z., Kortum, P., Xiong, J., \& Wallach, D. S. (2018). 2FA Might Be Secure, But It’s Not Usable: A Summative Usability Assessment of Google’s Two-factor Authentication (2FA) Methods. \textit{Proceedings of the Human Factors and Ergonomics Society Annual Meeting}, 62(1), 1141--1145.

\bibitem{ariyadi2025}
Ariyadi, A., \& Akbar, M. (2025). Public services in the digital age: A systematic review and bibliometric analysis of trends and challenges in digital governance. \textit{Tamalanrea: Journal of Government and Development (JGD)}, 2(2), 154--168.

\bibitem{bonneau2012}
Bonneau, J., Herley, C., van Oorschot, P. C., \& Stajano, F. (2012). The Quest to Replace Passwords: A Framework for Comparative Evaluation of Web Authentication Schemes. \textit{IEEE Symposium on Security and Privacy}.

\bibitem{bonneau2015}
Bonneau, J., Herley, C., van Oorschot, P. C., \& Stajano, F. (2015). Passwords and the evolution of imperfect authentication. \textit{Communications of the ACM}, 58(7), 78--87.

\bibitem{ciolino2019}
Ciolino, S., Parkin, S., \& Dunphy, P. (2019). Of two minds about two-factor: Understanding everyday FIDO U2F usability through device comparison and experience sampling. \textit{Proceedings of the 15th Symposium on Usable Privacy and Security (SOUPS 2019)}, 343--356.

\bibitem{colnago2018}
Colnago, J., Devlin, S., Oates, M., Swoopes, C., Bauer, L., Cranor, L., \& Christin, N. (2018). ``It’s not actually that horrible.'' \textit{Proceedings of the 2018 CHI Conference on Human Factors in Computing Systems}.

\bibitem{das2018}
Das, S., Dingman, A., \& Camp, L. J. (2018). Why Johnny Doesn’t Use Two Factor: A Two-Phase Usability Study of the FIDO U2F Security Key. \textit{Financial Cryptography and Data Security}, 160--179.

\bibitem{das2020}
Das, S., Wang, B., Kim, A., \& Camp, L. J. (2020). MFA is A Necessary Chore!: Exploring User Mental Models of Multi-Factor Authentication Technologies. \textit{HICSS}.

\bibitem{das2021}
Das, S., Kim, A., \& Camp, L. J. (2021). Organizational Security: Implementing a Risk-Reduction-Based Incentivization Model for MFA Adoption. \textit{SSRN Electronic Journal}.

\bibitem{debruijn2017}
de Bruijn, H., \& Janssen, M. (2017). Building Cybersecurity Awareness: The need for evidence-based framing strategies. \textit{Government Information Quarterly}, 34(1), 1--7.

\bibitem{dienlin2015}
Dienlin, T., \& Trepte, S. (2015). Is the privacy paradox a relic of the past? An in-depth analysis of privacy attitudes and privacy behaviors. \textit{European Journal of Social Psychology}, 45(3), 285--297.

\bibitem{eba2019}
European Banking Authority (EBA). (2019). \textit{Leitlinien zu IKT- und Sicherheitsrisikomanagement} (EBA/GL/2019/04).

\bibitem{enisa2024}
European Union Agency for Cybersecurity (ENISA). (2024). \textit{ENISA threat landscape 2024: July 2023 to June 2024}. Publications Office.

\bibitem{enisa2025}
European Union Agency for Cybersecurity (ENISA). (2025). \textit{Technical implementation guidance on cybersecurity risk management measures (NIS2)}.

\bibitem{enhuber2023}
Enhuber, M., Schwenk, J., \& Holz, T. (2023). Protection Motivation Theory in the Context of MFA: Usability versus Security. \textit{Proceedings of the 2023 ACM Workshop on Privacy and Electronic Society}, 198--208.

\bibitem{eu2014}
Europäische Union. (2014). Verordnung (EU) Nr. 910/2014 über elektronische Identifizierung und Vertrauensdienste (eIDAS). \textit{Amtsblatt der Europäischen Union}, L 257/73.

\bibitem{eu2015}
Europäische Union. (2015). Richtlinie (EU) 2015/2366 über Zahlungsdienste im Binnenmarkt (PSD2). \textit{Amtsblatt der Europäischen Union}.

\bibitem{eu2016}
Europäische Union. (2016). Verordnung (EU) 2016/679 (Datenschutz-Grundverordnung). \textit{Amtsblatt der Europäischen Union}, L 119/1.

\bibitem{eu2022a}
Europäische Union. (2022a). Richtlinie (EU) 2022/2555 (NIS2-Richtlinie). \textit{Amtsblatt der Europäischen Union}.

\bibitem{eu2022b}
Europäische Union. (2022b). Verordnung (EU) 2022/2554 (DORA). \textit{Amtsblatt der Europäischen Union}.

\bibitem{fagan2016}
Fagan, M., \& Khan, M. M. H. (2016). Why do they do what they do?: A study of what motivates users to (not) follow computer security advice. \textit{SOUPS 2016}, 59--75.

\bibitem{florencio2007}
Florêncio, D., \& Herley, C. (2007). A large-scale study of web password habits. \textit{Proceedings of the 16th International Conference on World Wide Web}, 657--666.

\bibitem{govindraj2024}
Govindraj, C.V. (2024). Digital Governance and Cybersecurity: Challenges in the Age of E-Governance. \textit{ShodhKosh}, 5(1), 1777--1781.

\bibitem{gratian2018}
Gratian, M., Bandi, S., Cukier, M., Dykstra, J., \& Ginther, A. (2018). Correlating human traits and cyber security behavior intentions. \textit{Computers \& Security}, 73, 345--358.

\bibitem{hu2014}
Hu, V. C., Ferraiolo, D., Kuhn, R., Schnitzer, A., Sandlin, K., Miller, R., \& Scarfone, K. (2014). \textit{Guide to Attribute Based Access Control (ABAC) Definition and Considerations}. NIST Special Publication 800-162.

\bibitem{keil2024}
Keil, M., \& Zugenmaier, A. (2024). Evaluating the evaluation criteria for account-recovery procedures in passwordless authentication. \textit{Gesellschaft für Informatik (GI)}, 73--84.

\bibitem{kim2011}
Kim, D., \& Solomon, M. G. (2011). \textit{Fundamentals of Information Systems Security}. Jones \& Bartlett Learning.

\bibitem{krol2015}
Krol, K., Philippou, E., De Cristofaro, E., \& Sasse, M. A. (2015). ``They brought in the horrible key ring thing!'' Analysing the Usability of Two-Factor Authentication in UK Online Banking. \textit{USEC}.

\bibitem{krykavska2023}
Krykavska, I. (2023). Regulatory and institutional support of digitalization and cybersecurity of the public administration system in the EU. \textit{Bulletin of Lviv Polytechnic National University}, 10(37), 148--152.

\bibitem{kunke2021}
Kunke, J., Wiefling, S., Ullmann, M., \& Lo Iacono, L. (2021). Evaluation of Account Recovery Strategies with FIDO2-based Passwordless Authentication. \textit{ArXiv}.

\bibitem{lang2016}
Lang, J., Czeskis, A., Balfanz, D., Schilder, M., \& Srinivas, S. (2016). Security Keys: Practical Cryptographic Second Factors for the Modern Web. \textit{Springer, Berlin, Heidelberg}.

\bibitem{lassak2024}
Lassak, L., Markert, P., Golla, M., Stobert, E., \& Dürmuth, M. (2024). A comparative long-term study of fallback authentication schemes. \textit{CHI '24}, 1--19.

\bibitem{lyastani2020}
Lyastani, S. G., Schilling, M., Neumayr, M., Backes, M., & Bugiel, S. (2020). Is FIDO2 the Kingslayer of User Authentication? A Comparative Usability Study of FIDO2 Passwordless Authentication. \textit{IEEE Symposium on Security and Privacy (SP)}.

\bibitem{mraihi2011}
M’Raihi, D., Machani, S., Pei, M., \& Rydell, J. (2011). \textit{TOTP: Time-Based One-Time Password Algorithm}. RFC 6238.

\bibitem{malhotra2004}
Malhotra, N. K., Kim, S. S., \& Agarwal, J. (2004). Internet Users’ Information Privacy Concerns (IUIPC): The Construct, the Scale, and a Causal Model. \textit{Information Systems Research}, 15(4), 336--355.

\bibitem{mark2022}
Mark, T., Smith, J., \& Johnson, R. (2022). Spillover effects of information security policies between work and home domains. \textit{Journal of Information Security and Applications}, 65, 42--51.

\bibitem{marky2022}
Marky, K., Ragozin, K., Chernyshov, G., Matviienko, A., Schmitz, M., Mühlhäuser, M., Eghtebas, C., \& Kunze, K. (2022). ``Nah, it’s just annoying!'' A Deep Dive into User Perceptions of Two-Factor Authentication. \textit{ACM Transactions on Computer-Human Interaction}, 29(5).

\bibitem{miller2012}
Miller, K. W., Voas, J., \& Hurlburt, G. F. (2012). BYOD: Security and Privacy Considerations. \textit{IT Professional}, 14(5), 53--55.

\bibitem{ogorman2003}
O’Gorman, L. (2003). Comparing passwords, tokens, and biometrics for user authentication. \textit{Proceedings of the IEEE}, 91(12), 2021--2040.

\bibitem{ometov2018}
Ometov, A., Bezzateev, S., Mäkitalo, N., Andreev, S., Mikkonen, T., \& Koucheryavy, Y. (2018). Multi-Factor Authentication: A Survey. \textit{Cryptography}, 2(1), 1--31.

\bibitem{otta2023}
Otta, S. P., Panda, S., Gupta, M., \& Hota, C. (2023). A systematic survey of multi-factor authentication for cloud infrastructure. \textit{Future Internet}, 15(4), 146.

\bibitem{owasp2025}
OWASP. (2025). \textit{Multifactor authentication - OWASP cheat sheet series}.

\bibitem{pearman2019}
Pearman, S., Thomas, J., Naechni, P., Habib, H., Ela, L., \& Egelman, S. (2019). Why People (Don't) Use Password Managers. \textit{SOUPS 2019}, 319--338.

\bibitem{reese2019}
Reese, K., Smith, T., Dutson, J., Armknecht, J., Cameron, J., \& Seamons, K. (2019). A usability study of five two-factor authentication methods. \textit{SOUPS 2019}, 357--370.

\bibitem{reynolds2018}
Reynolds, J., Smith, T., Reese, K., Dickinson, L., Ruoti, S., \& Seamons, K. (2018). A Tale of Two Studies: The Best and Worst of YubiKey Usability. \textit{IEEE Symposium on Security and Privacy}, 872--888.

\bibitem{staines1980}
Staines, G. L. (1980). Spillover versus compensation: A review of the literature on the relationship between work and nonwork. \textit{Human Relations}, 33(2), 111--129.

\bibitem{temoshok2025a}
Temoshok, D., et al. (2025). \textit{Digital identity guidelines: Authentication and authenticator management (NIST SP 800-63B-4)}. National Institute of Standards and Technology.

\bibitem{temoshok2025b}
Temoshok, D., et al. (2025). \textit{Digital identity guidelines (NIST SP 800-63-4)}. National Institute of Standards and Technology.

\bibitem{ulqinaku2021}
Ulqinaku, E., Assal, H., Abdou, A., Chiasson, S., \& Capkun, S. (2021). Is real-time phishing eliminated with FIDO? \textit{30th USENIX Security Symposium}, 3811--3828.

\bibitem{vance2020}
Vance, A., Siponen, M., \& Pahnila, S. (2020). Examining the spillover effect of corporate information security awareness on personal security behaviors. \textit{Information Systems Journal}, 30(1), 120--145.

\bibitem{wibowo2022}
Wibowo, J., \& Ramli, K. (2022). Impact of implementation of information security risk management and security controls on cyber security maturity. \textit{Jurnal Sistem Informasi}, 18(2), 1--17.
    
\bibitem{wiefling2019}
Wiefling, S., Lo Iacono, L., \& Dürmuth, M. (2019). Is This Really You? An Empirical Study on Risk-Based Authentication Applied in the Wild. \textit{ICT Systems Security and Privacy Protection}, 134--148.

\bibitem{bortz2006}
Bortz, J., Döring, N. (2006). Quantitative Methoden der Datenerhebung. In: Forschungsmethoden und Evaluation. Springer-Lehrbuch. Springer, Berlin, Heidelberg. https://doi.org/10.1007/978-3-540-33306-7_4

\bibitem{bortz2016}
Döring, N., & Bortz, J. (2016). Forschungsmethoden und Evaluation in den Sozial- und Humanwissenschaften (5. Aufl.). Springer. https://doi.org/10.1007/978-3-642-41089-5





\end{thebibliography}

